\documentclass[12pt]{article}

\usepackage[a4paper, total={6in, 8in}]{geometry}
\usepackage{textcomp}

\begin{document}
\setlength{\parindent}{0pt}

\def \t {\textrightarrow}
\def \v {\vspace{1em}}
\def \bi {\begin{itemize}}
\def \ei {\end{itemize}}
\def \s[#1] {\section*{#1}}
\def \ss[#1] {\subsection*{#1}}
\def \sss[#1] {\subsubsection*{#1}}

\s[D'Annunzio]
Nel fare della vita un'opera d'arte si colloca anche l'espressione "il verso è tutto" \t pregnanza del verso, e gli anni assidui con ambizione alle stelle \\
Prima fase a roma, poi a napoli \t danno il ritratto di un uomo nel quale la vita e l'arte si congiungono \\
Per la sua vitalita ha intemperanze \t e aveva gia dato il meglio di se quando incontro Giseulda Zucconi \t istintivita erotica \\
Abbiamo testimonianze che arrivano da lettere con compagni di studio \t "La guerra come soleggiere del mondo" e "Il bagno di sangue" ???? di Papini \t tratteggia un d annunzio all apice della fame, e tratteggia personalita particolare \\
Nella sua presentazione Papini compone "D'annunzio" nella pubblicazione "Passato remoto" ?? \\
Il sensualismo d'annunziano, ovvero la percezione attraverso i 5 sensi \t ma la percezione per lui vuol dire anche altro \\
Papini e Prezzolini da ricordare, "La firenze delle riviste"

\v

La ragazza, chiamata da lui Lalla, inaugura la sua vita sentimentale tormentata e infuocata \t i suoi incontri sono molteplici, e incontra donne di ogni tipo \\
D'annunzio è imprenditore di se stesso \\
La tendenza ad autopromuoversi \\
Maria Hardouin \t duchessa di Galles, nel 83 \t gia si rende protagonista di una fuga d'amore \t questa fuga d'amore desta scalpore, e il matrimonio con lei sara fuoco e fiamme \t si stanchera presto, anche se avranno 3 figli \\
Lui la lascia \\
Matrimonio tra arte e vita \t quando pero a napoli condurra la sua vita infuocata, incontrera Maria Gravina con cui avra una figlia \t lui la tradira \\
Gli anni ??? sono gli anni in cui la voce di Nietzsche arriveranno nella sua vita \t nella sua filosofia, lui trovera un senso allo suo stile di vita, che è in sintonia con il pensiero di Nietzsche \t "Il trionfo della morte" del 94 e 95 "Le vergini delle rocce", sono le produzioni associate al fuoco nietzschiano \\
D'annunzio pero senza la passione forte non sarebbe nulla \t passione forte porta a Alesiana Due ??????? \t attrice che conosce d'annunzio e che per lungo periodo sembra far vivere a d annunzio un amore totalizzante (ma lo stesso non fedele) \\
Se il Vittoriale è la reggia di ??, la Capponcina è un altra reggia del tum tum tum \t Capponcina cambia la location e siamo nel periodo nietzschiano, nel 95, e tra d'Annunzio e Dure ?? si crea relazione \\
Sono due ville che sono colli d'amore dove i due si frequentano e stanno \\
Vittoriale \t in provicia di Brescia, e Capponcina \t colli di ?? \\
Anche la capponcina presenta le caratteristiche del vivere inimitabile di d annunzio

\v

Con d'annunzio poeta vi è il "Libro delle laudi" \t Asterope, Merope e Alcyone \t la + imp. è Alcyone è importante \\
Trionfo del sensualismo \t la manifestazione del poeta presenta una riconferma di quanto detto: sensualismo, passione per mimesis uomo-natura, e la cosi detta naturalizzazione dell uomo e umanizzazione della natura (come nella pioggia del pineto e la notte fiesolana) \\
Nel frattempo frequenta anche delle donne a parigi, dalle quali ottiene passione folle \t "La citta morta" \t e tradisce l'Aduse \\
"La figlia di orio" \t opera teatrale di successo \\
Nel 1904 aduse si stanca e la loro relazione si sfalda \t lui abbandona la capponcina, i creditori lo assaltano e il poeta è desolato \t arriva a parigi, dove si "ritira" \\
In questo periodo pubblica nella rivista "Grande revue" (grande rivista) \t qua escono a puntate i capitoli del romanzo "Forse che si, forse che no"

\v

Ci si interroga in che rapporti sia stato annunzio con mussolini \t d'annunzio era entrato in politica tra le file della destra gia nel 97 \\
Poi sposta le sue preferenze verso l'estrema sinistra \t un volteggiare politicamente \t lo schieramento politico si sceglieva anche sull'onda di una vertenza politica \t ci si poteva trovare d'accordo con politiche di fazioni diverse \\
D'annunzio viene visto quindi come il volta faccia \t d'annunzio andava pero solo verso la vita, che viene distrutta dalla sua mania di accumulare, di giocare, dal voyerismo (revenge porn in qualche modo) \\
Visione di panismo d'annuzniano \t immersione totale nella natura \\
Mussolini, se da un lato coglie in d'annunzio un ottimo propagandista, dall'altro ne è intimidito \t d'annunzio scrive anche i discorsi per mussolini, e ha patrocinato un efficacia assertiva e di potere da parte di Mussolini \\
Si discute per l interventismo o neutralismo nella guerra \\
Nel 1915 torna con lo scoppio della guerra in italia \t e qui si capisce come d'annunzio si sente a rappresentare la penisola: si sente il vate della patria \\
Per lui la guerra è occasione per ergersi a vate della patria \t si vede come eroe di guerra nelle sue opere \t scrive lettere ad intellettuali, in cui si mostra come eroe di guerra (citata lettera, chiedere) \\
Memento audere semper (ricordati di usare sempre) \t espressione dell'audacia d'annunziana \t il vitalismo d'annunziano \\
Durante l'impero come fa l'oratoria ed essere autonoma? si deve incardinare nel potere \t anni in cui si diffono i mezzi di comunicazione, come la radio \t non c'è lentezza nella comunicazione \\
Volonta di sfidare il vuoto \t ovvero librarsi in un impresa enorme per l epoca \t Bava Beccaris \t ovvero quello si sorvolare su Vienna, che lui fa = messaggio di forte fiducia nella tecnologia, che riporta al futurismo \\
Il mondo sta cambiando, e l uomo vuole sfidare se stesso \\
A Maffeo Sciarra in una lettera, si parla di un esteta squattrinato di d annunzio \t sempre propenso a mettersi in mostra \\
A un certo punto Mussolini lo relega a un ruolo secondario \t come mai, se incarnava lo spirito nazionalistico? mussolini sia stato diffidente probabilmente di d annunzio, perche era troppo elevato nella sua caricatura intellettuale \\
Cosi si staccano, e lui va nell'esilio d'orato del Vittoriale sul lago di garda \t dimora spettacolare in Gardone Riviera
\end{document}
